\section*{Chapter \ref{ch:b2:expectations-and-lucas-critique} --- Expectations and the Lucas Critique}

\begin{solution}{exo:b2:expectations-and-lucas-critique:1}
% figdata: cbi, cbi_fit
Highest Spain, 9.95 per cent; lowest Germany, 3.19 per cent (Switzerland 3.42). Across these countries, each extra point on the
independence index goes with average inflation 1.83 percentage points lower.
\end{solution}

\begin{solution}{exo:b2:expectations-and-lucas-critique:2}
% figdata: phillips_dec
1960s: inflation 2.55 per cent and unemployment 4.78 on average, along a downward-sloping curve. 1970s: both higher, 7.22 and 6.22.
1980s: inflation down to 4.91, unemployment up to 7.28. The curve shifted up in the 1970s and back down in the 1980s.
\end{solution}

\begin{solution}{exo:b2:expectations-and-lucas-critique:3}
% figdata: it_summary
8.56 per cent five years before, 6.37 at adoption, 3.03 ten years after. The chart has no comparison group: inflation was
already falling before adoption and fell in many countries that did not adopt targets, and countries may have adopted
targets when they expected inflation to fall.
\end{solution}

\begin{solution}{exo:b2:expectations-and-lucas-critique:4}
The forecast $\pi_{t-1}$ falls short by one point every year: the error is $+1$ five years running. \omterm{def:b2:expectations-and-lucas-critique:rational-expectations}{Rational expectations} require
errors of mean zero that cannot be predicted from what was known; an error of the same sign year after year is
predictable, and a forecaster who noticed it would add a point.
\end{solution}

\begin{solution}{exo:b2:expectations-and-lucas-critique:5}
% code: exercises/solutions.py bias, expected_loss
Discretion: $\pi = 2 + 1 \times 0.5 \times 2 = 3$, $u = 5$ ($u^*$), loss $\tfrac12 \times 1^2 + \tfrac12 \times 2^2 = 2.5$. Commitment: $\pi = 2$, $u = u^*$, loss $\tfrac12 \times 2^2 = 2$. With $k = 1$ the bias is $1 \times 0.5 \times 1 = 0.5$.
\end{solution}

\begin{solution}{exo:b2:expectations-and-lucas-critique:6}
% code: listings/lucas.py
The intercept is $u^* + a\pi^{\mathrm T}$, so it rises by $a \times 2 = 4$. The listing's estimates rise from 8.92 to 13.25, by 4.33: close, the difference
being sampling noise in forty years of data.
\end{solution}

\begin{solution}{exo:b2:expectations-and-lucas-critique:7}
$E_{t-1}\pi_t = 2 - 0.5 \times (7 - 5) = 1$ per cent. With $\varepsilon_t = 1$, $\pi_t = 2$ and $u_t = 5 - 2 \times (2 - 1) = 3$. The systematic part of the rule (the response to $u_{t-1}$) moves expected and actual
inflation together and has no effect on $u_t$; only the surprise does.
\end{solution}

\begin{solution}{exo:b2:expectations-and-lucas-critique:8}
Under the 2 per cent target, $u_t = 5 - 2(\pi_t - \pi_{t-1})$ and, with surprises independent from year to year, the regression of $u$ on $\pi$ still has slope about $-2$ and intercept
about $5 + 2 \times 2 = 9$, so it again promises unemployment of about 1 per cent at 4 per cent inflation. This time it is right for one
year: in the first year of the new target inflation rises by about 2 points more than expected, and unemployment falls to
about 1 per cent. From the second year expectations have caught up, and unemployment returns to about 5 per cent. Keeping it
low would need inflation that rises every year.
\end{solution}

\begin{solution}{exo:b2:expectations-and-lucas-critique:9}
% code: exercises/solutions.py decade_inflation
Medians across the eighteen countries: 1950s 3.8 per cent, 1960s 3.6, 1970s 9.0, 1980s 7.1, 1990s 2.4, 2000s 2.1, 2010s 1.3. The highest country average was
Portugal's in both the 1970s (18.9 per cent) and the 1980s (17.4).
\end{solution}

\begin{solution}{exo:b2:expectations-and-lucas-critique:10}
Substituting $u - u^* = -a(\pi - \pi^e)$, the bank minimises $\tfrac12(\pi - \pi^{\mathrm T})^2 + \tfrac12w(k - a(\pi - \pi^e))^2$. The derivative is $(\pi - \pi^{\mathrm T}) - wa(k - a(\pi - \pi^e)) = 0$, so $\pi(1 + wa^2) = \pi^{\mathrm T} + wa^2\pi^e + wak$, which is \eqref{eq:b2:expectations-and-lucas-critique:reaction}.
Setting $\pi^e = \pi$: $\pi - \pi^{\mathrm T} = wak$, and $u = u^*$. The derivative of the bias with respect to $a$ is $wk > 0$: when a surprise buys more unemployment, the
temptation is greater, and so is the inflation that removes it.
\end{solution}

\begin{solution}{exo:b2:expectations-and-lucas-critique:11}
Now $u - u^* + k = k + e - a(\pi - \pi^e)$, and the first-order condition gives $(\pi - \pi^{\mathrm T})(1 + wa^2) = wa^2(\pi^e - \pi^{\mathrm T}) + wa(k + e)$. Taking expectations, with $\pi^e = E\pi$: $\pi^e - \pi^{\mathrm T} = wak$. Substituting,
$(\pi - \pi^{\mathrm T})(1 + wa^2) = wak(1 + wa^2) + wae$, so $\pi = \pi^{\mathrm T} + wak + \dfrac{wa}{1 + wa^2}e$. Then $\pi - \pi^e = \dfrac{wa}{1 + wa^2}e$ and $u = u^* + e - \dfrac{wa^2}{1 + wa^2}e = u^* + \dfrac{1}{1 + wa^2}e$.
\end{solution}

\begin{solution}{exo:b2:expectations-and-lucas-critique:12}
With the rule known, $\pi^e = \pi^{\mathrm T}$ and $u = u^* + (1 - ac)e$. Society's expected loss is $\tfrac12c^2\sigma^2 + \tfrac12\lambda(k^2 + (1 - ac)^2\sigma^2)$. The derivative in $c$, $c\sigma^2 - \lambda a(1 - ac)\sigma^2 = 0$, gives $c = \lambda a/(1 + \lambda a^2)$, the
same response as discretion with $w = \lambda$. The two policies therefore differ only in the bias, and the rule's loss is lower by $\tfrac12(\lambda ak)^2$.
\end{solution}

\begin{solution}{pb:b2:expectations-and-lucas-critique:1}
% code: exercises/solutions.py BRENLAND, bias, responses, expected_loss, commitment_loss, best_weight
\textbf{Part I.}
1. $u^* = 5$, $a = 0.5$, $\pi^{\mathrm T} = 2$, $k = 5 - 3 = 2$, $\lambda = 1$.
2. That the natural rate is higher than society would like, for instance because of taxes or market power; policy would
like unemployment at 3 per cent.
3. Small deviations cost little and large ones a lot, which gives the bank a reason to spread deviations across inflation
and unemployment.
4. $\tfrac12 \times 0 + \tfrac12 \times (5 - 3)^2 = 2$.
5. No: with \omterm{def:b2:expectations-and-lucas-critique:rational-expectations}{rational expectations} unemployment is $u^*$ on average whatever the policy, so the best is to add no inflation; the loss
2 is the minimum.
\textbf{Part II.}
6. $\pi = (2 + 0.25\pi^e + 1)/1.25 = 2.4 + 0.2\pi^e$.
7. $\pi = 2.4 + 0.2\pi$, so $\pi = 3$.
8. $u = 5$.
9. $\tfrac12 \times 1 + \tfrac12 \times 4 = 2.5$.
10. 2.0.
11. If people expected 2, the bank's best choice would be $2.4 + 0.4 = 2.8$, not 2: the announcement is not in its interest once expectations are
set, and people know it.
\textbf{Part III.}
12. $w \times 0.5 \times 2 = w$; 0.5 for $w = 0.5$.
13. Inflation $\dfrac{0.5w}{1 + 0.25w}e$ and unemployment $\dfrac{1}{1 + 0.25w}e$: 0.4 and 0.8 times $e$ for $w = 1$; 0 and $1 \times e$ for $w = 0$.
14. $\tfrac12(1 + 0.16 \times 4) + \tfrac12(4 + 0.64 \times 4) = 0.82 + 3.28 = 4.10$.
15. $\tfrac12 \times 0.64 + 3.28 = 3.60$.
16. $\tfrac12(0 + 0) + \tfrac12(4 + 4) = 4.00$: lower than with $w = 1$; removing the bias is worth more than the lost stabilisation here.
17. Bias 0.5, responses 0.222 and 0.889: loss 3.804.
\textbf{Part IV.}
18. $\tfrac12\bigl(w^2 + 4\bigl(\tfrac{0.5w}{1 + 0.25w}\bigr)^2\bigr) + \tfrac12\bigl(4 + 4\bigl(\tfrac{1}{1 + 0.25w}\bigr)^2\bigr)$.
19. $w = 0.425$.
20. 3.798, which is 0.198 above the best commitment (3.60) and 0.302 below discretion with $w = 1$.
21. Bias 0.42; responses 0.19 for inflation and 0.90 for unemployment.
22. At $w = 0$ a small rise in $w$ improves stabilisation at first order and adds a bias whose cost, $\tfrac12w^2$, rises only at second order; at
$w = 1$ the responses are the best ones, so a small cut in $w$ costs stabilisation only at second order and lowers the bias cost at
first order.
23. With $\sigma^2 = 16$ it rises to 0.71: stabilisation matters more. With $k = 4$ it falls to 0.18: the bias matters more.
24. Under discretion inflation exceeds the target by $wak$, here 1 point with $w = 1$, with no gain in unemployment; the expected social loss is
smallest with a bank that weighs unemployment at $w = 0.425$, between a pure inflation targeter and society's own weight.
\end{solution}
